\section{Introduction and the approximation choices}
\label{sec:intro}
The useful question for a finite asymptotic formula is what its
retained terms buy. Here there are two separate resolutions: the
curvature of the arithmetic carriers and the order of their
Riemann--Siegel completion. Improving either in isolation need not
improve the approximation to the canonical field. This paper derives
their matching, then controls the constants sufficiently well to let
the chosen order increase with height.

The development synthesizes the comparison and evaluation results in
recurrence Notes 02--06, 09--11, 16 and 18
\cite{rec02,rec03,rec04,rec05,rec06,rec09,rec10,rec11,rec16,rec18}. The comparison is always to the
same complete theta-source heat flow. In particular, a finite
polynomial hierarchy is not a truncation of the theta source or
of the retained arithmetic population. Its exact heat evolution
is useful structure, but the error is proved by matching initial
data, local remainders, and tails.

\subsection{What survives from the preceding ansatz investigation}
\label{subsec:retention}
The preceding investigation separated the elementary outer carrier
from the complete short arithmetic approximation
\cite{ans03,ans04,ans05,ans10,ans12}. We retain that distinction.
Put \(L(z)=\Log(z/(4\pi))\), with the branch continued from the
positive real axis, and define
\begin{subequations}\label{eq:outer-arms}
\begin{align}
 A_{\mathrm o}(t,z)&=\frac{\pi^{1/4}}{16}z^{7/4}
 \exp\left\{-\frac{\pi z}{8}
       +\frac t{16}\left(L(z)^2-\frac{\pi^2}{4}\right)\right\},\\
 \Theta_{\mathrm o}(t,z)&=\frac z4(L(z)-1)
                  +\frac{\pi t}{16}L(z)+\frac{7\pi}{8}.
\end{align}
\end{subequations}
The elementary cosine \(A_{\mathrm o}\cos\Theta_{\mathrm o}\)
describes the first carrier. Its adequacy is regime-dependent.
In the negative-time regime studied here, the full arithmetic
population cannot be replaced by that carrier using the present
estimates.

For a fixed chart cutoff \(N=\lfloor\sqrt{X/(4\pi)}\rfloor\), the
earlier simplest complete short formula is
\begin{subequations}\label{eq:short}
\begin{align}
 F_{\mathrm{short}}(t,z)
 &=A_{\mathrm o}(t,z)\sum_{n=1}^{N}w_n(t,z)\cos\phi_n(t,z)
       +\e^{-\pi^2t/64}\mathcal R_N(z),\\
 w_n(t,z)&=n^{-1/2}
     \exp\left\{\frac t4\bigl((\log n)^2-L(z)\log n\bigr)\right\},\\
 \phi_n(t,z)&=\Theta_{\mathrm o}(t,z)
                  -\left(\frac z2+\frac{\pi t}{8}\right)\log n,\\
 \mathcal R_N(z)&=(-1)^NG(z)a(z)^{-1/2}\mathcal C(a(z)-N),\\
 \mathcal C(v)&=\frac{\cos(2\pi(v^2-v)-\pi/8)}{\cos(2\pi v)},
 \qquad a(z)=\sqrt{z/(4\pi)}.
\end{align}
\end{subequations}
Here \(G\) is defined in \cref{eq:gamma-scale}, and removable
singularities in \(\mathcal C\) are filled analytically.
The cutoff is frozen when differentiating. Ansatz Note 05 proves
natural-scale \(O(X^{-3/4})\) field, spatial-through-three, and
first-time errors on \(1\le T=-t\le2\).
Ansatz Note 04 separately proves its signed crossover estimates on
\(\abs t\log(X/(4\pi))\le8\). These are different domains and
different weighted statements.

The first higher approximation replaces the elementary arithmetic
carriers by exact gamma modes with a curvature correction. Its
Riemann--Siegel completion must also improve: the earlier
target-dependent construction retains \(C_0,C_1\) in a bounded
real-data heat integral and has \(O(X^{-5/4})\) error. The present
frozen first level retains \(C_0,C_1,C_2\) and has
\(O(X^{-7/4})\) error. The source of this difference will be
derived in \cref{subsec:matching,subsec:low-orders}.

\begin{center}
\small
\begin{tabular}{@{}>{\raggedright\arraybackslash}p{.21\textwidth}>{\raggedright\arraybackslash}p{.29\textwidth}>{\raggedright\arraybackslash}p{.24\textwidth}>{\raggedright\arraybackslash}p{.14\textwidth}@{}}
\toprule
Approximation & Arithmetic terms retained & Completion retained &
Natural-scale error\\
\midrule
Earlier short formula &
All \(n\le N\), both elementary arms &
Leading \(\mathcal R_N\), with its heat factor &
\(X^{-3/4}\)\\[3pt]
Earlier first curvature correction &
All exact gamma modes, target-dependent quadratic correction &
Bounded integral through \(C_1\), exact transfer &
\(X^{-5/4}\)\\[3pt]
Frozen first level &
\(U_0+U_1\), all modes &
Bounded integral through \(C_2\) &
\(X^{-7/4}\)\\[3pt]
Frozen second level &
\(U_0+U_1+U_2\), all modes &
Bounded integral through \(C_4\) &
\(X^{-11/4}\)\\[3pt]
Frozen order \(p\ge1\) &
\(\sum_{j=0}^pU_j\), all modes &
Bounded integral through \(C_{2p}\) &
\(X^{-(p+3/4)}\)\\
\bottomrule
\end{tabular}
\end{center}
All rows apply on \(1\le T\le2\) with fixed buffers and sufficiently
large height; the frozen fixed-order rows extend to any declared
fixed compact interval \(0<\tau\le T\le B\).
Every integral row includes the exact change between the natural
cutoff at the integration point and the chart cutoff \(N\).
The table compares proved upper allowances, not pointwise
monotonicity of approximation error.

\subsection{Main conclusions and proof architecture}
The paper develops three connected conclusions. First,
\cref{thm:fixed,thm:growing} quantify the gain from each retained
fluctuation level, first at fixed order and then at controlled growing
order. The completion order must track the arithmetic curvature order:
improving one without the other can leave the error rate unchanged.

Second, \cref{thm:effective} turns this comparison into a terminating
finite evaluation construction. Analytic coefficient bounds, quadrature,
input sensitivity, and roundoff are allocated separately. Shifted gamma
expansions, exact Gauss root counts, pole-safe coefficient formulas, and
interior branch guards then supply every data contract. The accepted
input class consists of dyadic centers with strict certified descriptors
and exact rational queries; its separation dependence is explicit.

Third, \cref{thm:linear-onset,thm:bit-cost} make the construction
effective in height and cost. With \(\nu=\lceil r+b+2\rceil+1\),
a sufficient log-height is \(L_{\rm new}=64\nu\), improving the earlier
\(L_*=2^{22}\nu^4\) by the factor \(2^{16}\nu^3\). The bit-work bound
retains all \(N\asymp\sqrt X\) modes, while streaming avoids storing
them all. These conservative bounds establish a finite procedure; their
size does not assert practical evaluation at the sufficient onset.

The proofs proceed from the real initial-value heat formula to the local
generating function. Integrating its weighted remainder against Gaussian
moments preserves the useful powers of height; the full population costs
\(X^{1/4}\), and matching the initial Riemann--Siegel remainder fixes
\(K=2p\). True and artificial tails are kept separate. Only after their
order dependence is controlled do we change to the smaller outer envelope.
The evaluation budget is then proved as a modular theorem, supplied by
finite algorithms, and closed with explicit uniform calibration.

The analytical claims are proved from the named classical inputs.
Their source notes have same-model internal assessments, and this
manuscript provides an internal synthesis check. Supporting component
computations are linked in \cref{sec:effective}; they are not premises of
the theorems. Independent specialist validation, including the source
proofs behind the imported coefficient identities, remains open.
No claim of priority, infinite-series convergence, or a conclusion about
zeros from approximation accuracy alone is made.
